\section{Objectifs}
\begin{itemize}
\item Comprendre pourquoi les logarithmes ont simplifié les calculs.
\item Relier logarithme et exponentielle.
\item Utiliser les propriétés algébriques et analytiques de $\ln$.
\end{itemize}

\section{1. Idée historique}
Avant les calculatrices, transformer un produit en somme était un gain majeur.

Le logarithme vérifie :
\[
\ln(ab)=\ln a+\ln b.
\]

Les tables de logarithmes permettaient donc de remplacer des multiplications difficiles par des additions.

\section{2. Fonction réciproque}
Pour $x>0$ :
\[
\ln x=y\iff e^y=x.
\]

Donc :
\[
\ln(e^x)=x,\qquad e^{\ln x}=x.
\]

\section{3. Propriétés}
Pour $a,b>0$ :
\[
\ln(ab)=\ln a+\ln b,
\]
\[
\ln(a/b)=\ln a-\ln b.
\]

\section{4. Dérivée}
\[
 (\ln x)'=\frac1x.
\]

La fonction est strictement croissante sur $]0,+\infty[$.

\section{5. Limites}
\[
\ln x\to-\infty\quad(x\to0^+),
\]
\[
\ln x\to+\infty\quad(x\to+\infty).
\]

\section{6. Équations}
\[
\ln x=a\iff x=e^a.
\]

Les logarithmes permettent aussi d'isoler un exposant dans une équation exponentielle.

\section{Erreurs fréquentes}
\begin{itemize}
\item Écrire $\ln(a+b)=\ln a+\ln b$.
\item Oublier $x>0$.
\item Confondre logarithme et puissance.
\end{itemize}

\section{À retenir}
Le logarithme est à la fois un outil historique de calcul et une fonction essentielle de l'analyse moderne.
